\subsection{Generalized Lotka-Volterra Model} To model the dynamics of community that affected by species interactions, ordinary differential equations are used. Generally, the model would be a equation system where each of the differential equation describes the change of the population size of a species. Generalized Lotka-Volterra model (Eqn.~\ref{eqn:L-V_equation}) is well used. The system is written \begin{equation}\label{eqn:L-V_equation} \frac{\mathrm{d} N_i}{\mathrm{d} t} = N_i \left( r_i + \sum_{j=1}^{S} A_{ij} N_j \right), \end{equation} where there are $S$ species, $N_i$ is the population size of the $i$th species, $r_i$ represent the intrinsic rate of $i$th species, and $A_{ij}$ is the interaction strength of species $j$ on $i$, which is a real value. There are other models which include more effects. But GLV captured the environment capacity and species interaction, and is easy to deal with at the mean time. In this model, increase rate per population size is a linear term of all other population size, makes a simple form of Jacobian matrix around fixed point. Thus, the stability and stable point would be limited by two condition. Firstly, the feasible condition \begin{equation} \boldsymbol{N}^* = -\boldsymbol{A}^{-1}\boldsymbol{r} > 0 \end{equation} makes sure that there is a fixed point where all species coexist. Secondly, the stable condition \begin{equation*} \max_i \operatorname{Re} \left[ \lambda_i(\boldsymbol{M}) \right] < 0, \end{equation*} provide a criteria that the fixed point is stable. Only the parameter combination that fulfills both condition is viewed as a stable system.

\subsection{Diffusion-Reaction Model} Besides of interactions, spacial community dynamics must include dispersal of species. Dispersal is normally modelled as the spacial derivation of population size. As a result, models include dispersal is normally partial differential equations. The diffusion model, which is the second spacial derivative, is well-used and easy to derive. The diffusion model assume all individuals are dispersal randomly to both side. And also, it indicates that some individuals might spread very fast. In reality, dispersal of species would also be affected by the species interaction, geographical condition, and other random effect. However, it works well as a null model where we would assume that each individual spread symmetrically and independently. Other dispersal model will also be discussed. Based on which, the diffusion-reaction model is applied to simulate the spacial dynamic of biological system where the diffusion term represent random diffusion of biological molecular or individuals, as well as reaction term indicate the chemical reaction or interaction. In development biological and cellular biology, the boundary condition is normally set to be a periodical boundary, which indicate the space is topologically the same to a circle or a ball. For community ecology model, here we set a concrete boundary where flux in both side are set to $0$, to model the finite habitat of species. A general form of diffusion-reaction model is \begin{equation}\label{eqn:D-R model} \begin{aligned} \frac{\partial N_i}{\partial t} = D_{i}\frac{\partial^2 N_i}{\partial x^2} + f(N_j) \end{aligned} \end{equation} where $N_{i}$ is the amount of $i$th article (here, the population size), $D_{i}$ is the dispersal speed, and $f(N_j)$ is a function of all $N_{j}$, indicate the reaction between articles. Here we use the generalized Lotka-Volterra model as the reaction term. The space is set to $[-1,1]$. One species and two species cases are detailed discussed, and a special case of multiple species model is used to show more complex cases.

\subsubsection{One Species Invading Process} For one species, there is only one partial differential equation in the system, \begin{equation}\label{eqn:single_species} \frac{\partial u}{\partial t} = D \frac{\partial ^2 u} {\partial x^2} + f(u), \end{equation} and we set the initial condition as $u(x) = 1$ if $x<0$ and $u(x) = 0$ if $x>0$. The initial condition means that the population occupied the left part and invading the right part. If $f(u) = u(1-u)$, the model with logistic reaction term is called the Fisher-KPP model. $f(u)$ has two root $0$ and $1$. For the reaction dynamics, $u=0$ is a instable stationary point as well as $u=1$ is stable. As a result, the front connecting the occupied region where $u = 1$ and the empty space $u=0$ would corresponding to a trajectory connecting these two stationary point. Since $u=0$ is instable, the front would move to the instable side. By assuming a moving wave solution that fit the guess, we could set $z = x-ct$ such that $u(z) = u(x-ct)$. Thus the shape of the front is limited by a second ordered ODE \begin{equation*} \frac{\mathrm{d}^2 u}{\mathrm{d} z^2} + c \frac{\mathrm{d} u}{\mathrm{d} z} + u(1-u) = 0. \end{equation*} To make sure the existence of the solution, moving speed $c$ must large than $2$. The solution of $u(z)$ could only given in the form of a series, but the first term is already a good approximation which is a logistic function \begin{equation} u(z) = \frac{1}{1 + e^{z/c}} + \mathcal{O}(c^{-2}) \end{equation} To have a stationary front, where must be two stable state in both side of the space. A candidate model is the Allen-Chan model, where the reaction term is $f(u) = u(1-u)(u-a)$. The $(u-a)$ term could be understood as Allee effect which indicate the benefit of population size. This reaction function has two stable point $u = 0$ and $u = 1$, as well as an instable one $u = a$. As a result, the moving direction and speed would be determined by the potential of these two stable state. The potential $W$ is defined as $\frac{\mathrm{d} W(u)}{\mathrm{d} t} = f(u)$, then $W(0) = 0$, $W(1) = \frac{1}{12} - \frac{a}{6}$, moving speed is given by $c = \sqrt{2}(a - \frac{1}{2})$. The moving dynamics is fully determined by $a$, and the front kept stationary when $a = \frac{1}{2}$. The shape of moving front can be solved explicitly, \begin{equation} u(z) = \frac{1}{1 + e^{z/c}}, \end{equation} which is exactly the logistic function.

\subsubsection{Two Species Interaction-Dispersal Community} When inter-specific interaction is considered, two species model is the simplest to include it. A two species diffusion-reaction model with GLV reaction term is \begin{equation}\label{eqn:two_species} \begin{aligned} \frac{\partial N_1}{\mathrm{d}t} = D_{1}\frac{\partial^2 N_1}{\partial x^2} + N_1 \left( r_1 - \alpha_{1,1} N_1 - \alpha_{1,2} N_2 \right)\\ \frac{\partial N_2}{\mathrm{d}t} = D_{2}\frac{\partial^2 N_2}{\partial x^2} + N_2 \left( r_2 - \alpha_{2,2} N_2 - \alpha_{2,1} N_1 \right) \end{aligned}. \end{equation} To illustrate the coexistence, the system should be bistable. As a result, intrinsic rate is set to positive, and inter-specific competition is high, to make sure these two species cannot coexist locally. This system can be rescaled as \begin{equation} \begin{aligned} \frac{\partial u}{\mathrm{d}\tau} = \frac{\partial^2 u}{\partial z^2} + u \left( 1 - u - a_{1,2} v \right)\\ \frac{\partial v}{\mathrm{d}\tau} = D \frac{\partial^2 v}{\partial z^2} + r v \left( 1 - v - a_{2,1} u \right) \end{aligned}, \end{equation} where \begin{align*} u = \alpha_{1,1}/r_1 N_1 \quad v = \alpha_{2,2}/r_2 N_2 \quad z = \sqrt{r_1/D_1} x \quad \tau = r_1 t\\ r = r_2/r_1 \quad D = D_{2}/D_{1} \quad a_{1,2} = \frac{r_2 \alpha_{1,2}}{r_1 \alpha_{2,2}} \quad a_{2,1} = \frac{r_1\alpha_{2,1}}{r_2\alpha_{1,1}}. \end{align*} Four parameters are needed to describe the system, namely $D$, $r$, $a_{1,2}$, $a{2,1}$. We discuss the most simplest case where two species are symmetric, $D = 1$, $r = 1$, $a_{1,2} = a_{2,1}$. with proper approximation on the total population size, approximate solution could be given. For generalized model, we simulated the dynamics with different parameter combinations. We scanned dispersal rate $D$ ($0,2$) and intrinsic rate $r$ ($0,2$) and simulate the population distribution as, \begin{equation*}\label{eqn:balance_D-r} \begin{aligned} \partial_t N_1 & = D \cdot & \nabla^2 N_1 & + r \cdot &(1 - N_1 - 1.5 N_2)\\ \partial_t N_2 & = & \nabla^2 N_2 & + &(1 - N_2 - 1.5 N_1) \end{aligned}, \end{equation*} Similarly, we scanned dispersal rate $D$ in region $(0,2)$ as well as competing strength $a_{1,2}$ in interval $(0,2)$, with model, \begin{equation*}\label{eqn:balance_D-a} \begin{aligned} \partial_t N_1 & = D \cdot & \nabla^2 N_1 & + (1 - N_1 - a_{1,2} \cdot &1.5 N_2)\\ \partial_t N_2 & = & \nabla^2 N_2 & + (1 - N_1 - &1.5 N_2). \end{aligned} \end{equation*} Effect of intrinsic rate and competing strength are discussed separately, since these two parameter mainly related to the local dynamics as well as stability of GLV. The location $x^*$ where $u(x^*) = v(x^*)$ at the end of simulation is measured to show the moving speed of different condition as well as the coexistence criteria.

\subsubsection{Competing Multispecies Community} More generally, with multiple species, a GLV diffusion reaction model is \begin{equation}\label{eqn:multi_species} \begin{aligned} \frac{\partial N_i}{\partial t} = D_{i}\frac{\partial^2 N_i}{\partial x^2} + N_i \left( r_i - \sum_{j} A_{i,j} N_j \right) \end{aligned}, \end{equation} where $D_i$ is the diffusion rate of species $i$, $r_i$ is the intrinsic rate of that, and $A_{i,j}$ is the interaction strength between species $j$ and $i$. We focus on the bistable communities. As a result, the total community is instable, but there are two subgroup of the species combinations are stable, namely community $1$ and $2$. The dynamics could also be written in the following form, \begin{equation*} \begin{aligned} \partial N_{1,i} & = D_{1,i} \nabla^2 N_{1,i} + r_{1,i} N_{1,i} \left( 1 - \sum_{k} A_{ik} N_{1,k} - \sum_{\ell} A_{i\ell} N_{2,\ell} \right)\\ \partial N_{2,j} &= D_{2,j}\nabla^2 N_{2,j} + r_{2,j} N_{2,j} \left( 1 - \underbrace{\sum_{\ell} A_{j\ell} N_{2,\ell}}_{\text{intra-community}} - \underbrace{\sum_{k} A_{jk} N_{1,k}}_{\text{inter-community}} \right) \end{aligned}, \end{equation*} to divide the interaction term in GLV part as intra-community interactions and inter-community interactions. The following condition is discussed analytically, where diffusion rate is the same to all species, and species in the same stable subgroup are identical. More general case where diffusion rate are different or interaction strength is not the same are simulated to show the structure of multi-specific transition front.

\subsection{Simulation of the Partial Differential Equation Systems} For the nonlinear partial differential equations like the diffusion reaction model, it is normally very hard or sometimes impossible to get an exact solution. If it is a PDE system where there are more species, it get even harder. To show the dynamics as well as stability, we finite difference method to simulate the system, by discretize both space and time. We used the following two methods for rough simulation of parameter scanning as well as precise simulation of certain special cases. \subsubsection{Explicit Finite Difference Method} We first used the Euler method to discretize the PDE. The continuous space is divided by interval $\Delta x$, as well as time be divided by discrete time interval $\Delta t$. To simplify the equation, we write $N_{i}(x = k\Delta x, t = \ell \Delta t)$ as $u_{k,l}$. Then the spacial diffusion term is approximately replaced by the second-ordered difference term, \begin{equation*} \frac{\partial^2 N_i(k \Delta x,t)}{\partial x^2} \approx \frac{u_{k+1,\ell} - 2 u_{k,\ell} + u_{k-1,\ell}}{(\Delta x)^2} := \nabla^2 u_{k,\ell}. \end{equation*} The temporal term is approximately the difference of $u$ on time, \begin{equation*} \frac{\partial N_i(k, \ell \Delta t)}{\partial t} \approx \frac{u_{k,\ell+1} - u_{k,\ell}}{\Delta t} := \partial u_{k,\ell}. \end{equation*} Then the diffusion reaction system is discretized as \begin{equation} \partial u_{k,\ell} = D_{i} \nabla^2 u_{k,\ell} + f(\mathbf{u}_{k,\ell}), \end{equation} where $\mathbf{u}_{k,\ell}$ is the vector of $\{ u_{k,\ell} \}$ for all the $k$. Thus the relation between $\mathbf{u}_{k,\ell}$ and $\mathbf{u}_{k,\ell+1}$ is written, \begin{equation*} \frac{u_{k,\ell+1} - u_{k,\ell}}{\Delta t} = D_{i} \frac{u_{k+1,\ell} - 2 u_{k,\ell} + u_{k-1,\ell}}{(\Delta x)^2} + f(\mathbf{u}_{k,\ell}), \end{equation*} where $\mathbf{u}_{k,\ell+1}$ can be solved given $\mathbf{u}_{k,\ell}$. The explicit difference method is used to do rough simulation and scan the parameter space. \subsubsection{IMEX semi-implicit Method} Explicit finite difference method might be instable or divergence if there are rapid changes, especially when diffusion is faster than the interactions. In our simulation, we always assume that dispersal is at least a comparable process with local interactions, or the system would becomes several patches. As a result, we used the semi-implicit method to make the simulation more stable and precise. Similar to the explicit method, temporal differential term is the simple difference term of $u$. To make sure that the value of $u$ not be too diverge from the real value if the diffusion term change much faster, the function is changed to, \begin{equation} \partial u_{k,\ell} = D_{i} \nabla^2 u_{k,\ell+1} + f(\mathbf{u}_{k,\ell}), \end{equation} where the diffusion term is replace by difference approximation at $\ell+1$. To match the precise of the semi-implicit method, four-ordered difference is used to simulate the diffusion, \begin{equation*} \frac{\partial^2 N_i(k \Delta x,t)}{\partial x^2} \approx \frac{-u_{k+2,\ell} + 16 u_{k+1,\ell} - 30 u_{k,\ell} + 16 u_{k-1,\ell} - u_{k-2,\ell}}{12(\Delta x)^2}. \end{equation*} Thus the full simulation is, \begin{equation*} \frac{u_{k,\ell+1} - u_{k,\ell}}{\Delta t} = D_{i} \frac{-u_{k+3,\ell+1} + 16 u_{k+2,\ell+1} - 30 u_{k+1,\ell+1} + 16 u_{k,\ell+1} - u_{k-1,\ell+1}}{12(\Delta x)^2} + f(N_j(x,t)), \end{equation*} which is still an equation of $\mathbf{u}_{k,\ell}$ and $\mathbf{u}_{k,\ell+1}$, but not an explicit function. $\mathbf{u}_{k,\ell+1}$ should be solve from the implicit function. In each step, we use a sparse matrix of spacial difference to solve a linear function. Compared to the explicit one, this method consume more time to solve matrices. As a result, we only use it on the simulation of certain cases.

\subsection{Logistic Test for Equilibrium Distribution} The shape of the simulated population distribution is in a sigmoid-like form. Also, we know that the shape of moving wave front of the Allen-Chan model is logistic function. As a result, we try to test how well a logistic function could represent the stable population distribution, if there is species coexistence in the bistable community. For any species $N$, if the stable equilibrium abundance of a species in the first community is $K_1$, and in the second is $K_2$, a perfect logistic model is, \begin{equation} N_L(x) = \frac{K_1 - K_2}{1 + e^{-k(x-x_0)}} + K_2, \end{equation} where $x_0$ is the center of the front, as well as $k$ is the slope of the cline. When $x\rightarrow -\infty$, $N = K_2$, and when $x \rightarrow +\infty$, $N = K_1$. For the special case where species only exist in the first stable local community, the function is, \begin{equation*} N_L(x) = \frac{K}{1 + e^{-k(x-x_0)}}. \end{equation*} After the simulation, population distribution of each species is assumed to be converged. Then they are fit to a Logistic model as above. Fitting residuals are measured as $s(x) = N(x) - N_L(x)$. $R^2$ is computed for each regression.

\subsection{Measure $\beta$-Diversity} The measure of biodiversity is important for understanding the stability as well as distribution of species. $\alpha$-diversity represent the biodiversity in a well-mixed or nearly uniform community, which we called local community. $\gamma$-diversity, in the contrast, is the biodiversity of the community in a large region, called global community. There are many different measures of such biodiversity which depends on both the amount of species and the population size. We use the Shannon-Weaver index to measure the $\alpha$ and $\gamma$ diversity. The Shannon-Weaver index is defined form the perspective of information theory, measures the uncertainty of randomly observed species. For discrete samples, the measures are, \begin{align*}\label{eqn:alpha_gamma_dis} \alpha_k &= -\sum p_{k,i} \ln(p_{k,i}) \qquad p_{k,i} = \frac{N_{k,i}}{\sum_{j}N_{k,j}}\\ \gamma &= -\sum q_i \ln(q_i) \qquad q_i = \frac{\sum_{k}N_{k,i}}{\sum_{j}\sum_{k}N_{k,j}} \end{align*} Besides, $\beta$-diversity is a index to measure the change of community components. As a index showing the change of biodiversity along certain spacial condition, especially the change of species, it is also be viewed as measure of species turnover in space. Based on the above biodiversity measures, $\beta$ diversity is sometimes be defined as the fraction of $\gamma$ diversity and average $\alpha$ diversity. \begin{equation*}\label{eqn:beta_dis} \beta = \frac{\gamma}{\frac{1}{n} \sum_{k} \alpha_k} \end{equation*} If the distribution of species is known, that is, we have the continuous distribution $N_i(x)$, instead of samples $N_{k,i}$, the definition of different biodiversity measures can be extended with the distribution, for example, \begin{align*}\label{eqn:alpha_gamma_con} \alpha(x) &= -\sum p_{i}(x) \ln(p_{i}(x)) \qquad p_{i}(x) = \frac{N_{i}(x)}{\sum_{j}N_{j}(x)}\\ \gamma &= -\sum q_i \ln(q_i) \qquad q_i = \frac{\int_{x} N_{i}(x)}{\sum_{j} \int_{x} N_{j}(x)}, \end{align*} and thus \begin{equation*}\label{eqn:beta_con} \beta = \frac{\gamma}{\frac{1}{L} \int_{x} \alpha(x)}. \end{equation*} We applied this function on the simulated data as the real value of $\beta$-diversity. We then also derived the expression when assuming the species distribution is logistic.